\section{Исследование операций и задачи искусственного интеллекта. Экспертизы и неформальные процедуры. Автоматизация проектирования}

\subsection{Исследование операций}

\textbf{Исследование операций} --- прикладная математическая дисциплина, разрабатывающая методы обоснования решений при управлении сложными системами. Формальная задача обычно содержит множество допустимых решений $X$ и критерий эффективности:
\[
x^*=\arg\min_{x\in X}F(x)
\quad\text{или}\quad
x^*=\arg\max_{x\in X}F(x).
\]

К типичным задачам относятся распределение ресурсов, линейное и целочисленное программирование, транспортные и сетевые задачи, календарное планирование, управление запасами, массовое обслуживание и имитационное моделирование.

При нескольких критериях
\[
F(x)=(F_1(x),\ldots,F_m(x))
\]
решение $x_1$ \textbf{доминирует по Парето} решение $x_2$, если оно не хуже по всем критериям и строго лучше хотя бы по одному. Решения, не доминируемые другими допустимыми решениями, образуют множество Парето.

Для многоэтапных процессов важен \textbf{принцип оптимальности Беллмана}: хвост оптимальной стратегии, начинающийся из любого достигнутого состояния, сам должен быть оптимальным для соответствующей подзадачи.

\subsection{Задачи искусственного интеллекта}

\textbf{Искусственный интеллект} объединяет методы решения задач, требующих поиска, вывода, обучения или работы со знаниями. К классическим направлениям относятся:
\begin{itemize}
\item поиск в пространстве состояний и планирование;
\item представление знаний и логический вывод;
\item распознавание и классификация;
\item машинное обучение;
\item поддержка принятия решений.
\end{itemize}

Если задача представлена графом состояний, решение может искаться алгоритмами BFS/DFS или эвристическим поиском. Например, $A^*$ использует оценку
\[
f(n)=g(n)+h(n),
\]
где $g$ --- уже накопленная стоимость пути, $h$ --- эвристическая оценка оставшейся стоимости. Для поиска по дереву допустимость $h$ (эвристика не завышает истинную стоимость до цели) обеспечивает оптимальность $A^*$. Для поиска по графу с закрытым множеством состояний стандартным достаточным условием является также \emph{согласованность} (монотонность) эвристики. При только допустимой эвристике оптимальность сохраняется, если алгоритм допускает повторное открытие вершины при нахождении пути с меньшим $g$.

\subsection{Экспертные системы}

\textbf{Экспертная система} --- программа, оперирующая знаниями определённой предметной области для выработки рекомендаций или решения специализированных задач.

Классическая структура включает:
\begin{itemize}
\item базу знаний;
\item рабочую память фактов;
\item \textbf{машину вывода};
\item средства приобретения и редактирования знаний;
\item пользовательский интерфейс и, желательно, подсистему объяснения.
\end{itemize}

Знания могут представляться продукционными правилами, семантическими сетями, фреймами или логическими моделями. Продукционное правило имеет форму
\[
\text{ЕСЛИ условие, ТО действие/заключение}.
\]

При \textbf{прямом выводе} рассуждение начинается с известных фактов и применяет подходящие правила до получения заключения. При \textbf{обратном выводе} выбирается гипотеза и последовательно проверяются условия, достаточные для её подтверждения.

Разработка экспертной системы включает идентификацию задачи, концептуализацию предметной области, формализацию знаний, реализацию и тестирование. Трудность явного извлечения знаний у специалиста традиционно называют проблемой приобретения знаний.

\subsection{Экспертизы и неформальные процедуры}

Когда критерии трудно полностью формализовать или данных недостаточно, используют \textbf{экспертные оценки}. Типичные процедуры:
\begin{itemize}
\item ранжирование;
\item балльная оценка;
\item попарное сравнение;
\item интервью и анкетирование;
\item метод Делфи;
\item сценарный анализ и коллективное обсуждение.
\end{itemize}

Оценки нескольких экспертов могут агрегироваться, например
\[
R_i=\sum_{j=1}^{m}w_jr_{ij},
\qquad w_j\ge0,
\]
где $w_j$ характеризует вес/компетентность эксперта. Для групповой экспертизы важно проверять согласованность суждений, а не только усреднять числа.

Неформальные процедуры не противопоставляются математическим методам: часто они используются для выбора критериев, структуры модели или начальных альтернатив, после чего применяется формальная оптимизация.

\subsection{Автоматизация проектирования}

\textbf{Система автоматизированного проектирования (САПР)} --- организационно-техническая система, в которой проектные процедуры выполняются при взаимодействии человека и комплекса программно-технических средств.

В инженерной практике различают:
\begin{itemize}
\item \textbf{CAD} --- геометрическое и параметрическое проектирование;
\item \textbf{CAE} --- инженерный анализ и численное моделирование;
\item \textbf{CAM} --- технологическая подготовка производства.
\end{itemize}

При параметрическом проектировании объект описывается вектором параметров $x$. Тогда автоматизированный поиск варианта проекта естественно приводит к задаче
\[
F(x)\to\min,
\qquad g_i(x)\le0.
\]
Например, минимизировать массу конструкции при ограничениях по напряжениям и перемещениям.

Современный проектный цикл можно представить как
\[
\boxed{
\text{параметры}
\rightarrow
\text{CAD-модель}
\rightarrow
\text{CAE-расчёт}
\rightarrow
\text{оценка критериев}
\rightarrow
\text{оптимизация}
\rightarrow
\text{новые параметры}}.
\]
Методы ИИ и экспертные знания могут дополнять этот цикл при выборе вариантов и обработке трудноформализуемых требований, но результат должен проверяться инженерными ограничениями и расчётной моделью.

\subsection{Что важно сказать на экзамене}

Нужно показать связь четырёх частей билета: исследование операций формализует \textbf{выбор оптимального решения}; ИИ добавляет поиск, знания и обучение; экспертные процедуры работают с трудноформализуемой информацией; САПР переносит эти методы в автоматизированный проектный цикл.

\newpage
